\chapter{Referenz: Raytracer-Plugin}\label{kap:plugin-raytracer}
\index{Raytracer!Referenz}

Dieser Anhang listet alles auf, was das Plugin \plugin{Raytracer} anbietet. Die Bedienung mit
Beispielen beschreibt \autoref{kap:raytracer}.

\section{Steckbrief}\label{sec:ref-rt-steckbrief}

\begin{referenz}{@{}l >{\raggedright\arraybackslash}X@{}}
\toprule
Eigenschaft & Wert \\
\midrule\endhead
Kennung (\param{id}) & \code{raytracer} \\
Name & Raytracer \\
Version & 0.1.0 \\
Plugin-API & 1 \\
Einstiegspunkt & \code{raytracer\_plugin:RaytracerPlugin} \\
Beschreibung & Raytatouille: load optical systems, trace rays, irradiance on the image, optimise live. Needs the Python package raytatouille in the plugin environment. \\
Python-Anforderungen & \code{numpy}; \code{raytatouille>=0.9}; \code{rtt\_fusion>=0.2} (pip: \code{rtt-fusion}); \code{fieldtracing>=0.2} (pip: \code{photonpower-fieldtracing}) \\
Module & \datei{raytracer\_plugin.py} (Aktionen), \datei{rtt\_nurbs.py} (Umwandlung der Flächen in NURBS) \\
Beispiel & \datei{examples/singlet\_optim.rtt.json} \\
\bottomrule
\end{referenz}

\section{Aktionen}\label{sec:ref-rt-aktionen}

Alle Aktionen außer \gui{Section $\to$ group} erscheinen als Knopf im Panel \gui{Plugin}, im Menü
\menu{Plugins,Raytracer} und als Knopf der Gruppe \gui{Raytracer} in der Werkzeugspalte (dort
beschriftet mit den ersten drei Buchstaben). Die Reihenfolge ist die der Anwendung (alphabetisch
nach dem internen Namen). Jede Aktion ist ein Rückgängig-Schritt.

\begin{referenz}{@{}>{\raggedright\arraybackslash}p{3.0cm} >{\raggedright\arraybackslash}X >{\raggedright\arraybackslash}p{4.4cm}@{}}
\toprule
Aktion & Wirkung & Statusmeldung bei Erfolg \\
\midrule\endhead
\gui{Camera image} & Bestrahlungsstärke auf der letzten Fläche jedes gewählten Strahlengangs mit vielen Strahlen, als Bild im Fenster \gui{Image} und wahlweise als Feld auf der Fläche; Strahlengänge mit derselben letzten Fläche werden addiert & \texttt{camera image on }\textit{Fläche}\texttt{ (peak …) from … rays} \\
\gui{Follow the workflow} & ergänzt die Beschreibung der bearbeitbaren Daten der \code{rtt.*}-Blöcke im Workflow und zeichnet das System neu, wenn der Workflow es geändert hat. Wird von der Anwendung nach jeder Änderung des Node-Workflows still ausgeführt; schlägt nie fehl & leer \\
\gui{Load system} & liest die Systemdatei, lädt die Glaskataloge, zeichnet alle Flächen und Linsenränder in einer Gruppenstruktur nach dem Systembaum; ersetzt eine frühere Zeichnung & \textit{System}\texttt{: }\textit{n}\texttt{ surfaces} \\
\gui{New Raytatouille system} & legt im Node-Workflow ein neues System an (Quelle, Blende, Detektor, Zweig \code{main}) und zeichnet es; braucht RTT-Fusion mit \code{new\_system} & \texttt{new system \dq{}…\dq{}: source, stop, detector and the branch main – …} \\
\gui{Optimise} & Optimierer von Raytatouille mit Variablen und Gütefunktion des Systems; zeichnet nach jeder angenommenen Iteration neu; Kurve der Gütefunktion; \keys{Esc} endet mit dem letzten angenommenen Stand & \textit{Status}\texttt{ after }\textit{n}\texttt{ iterations, }$\varphi$\texttt{ = …} \\
\gui{Save system} & schreibt das aktuelle System (bzw. das System des Workflows) als \datei{.rtt.json} nach \gui{Save as} oder als \datei{\textit{System}.optimised.rtt.json} & \texttt{saved }\textit{Pfad} \\
\gui{Section $\to$ group} & nur im Kontextmenü eines Blocks \code{rtt.section} im Node-Workflow: zerlegt den Abschnitt in eine Gruppe mit einem Block je Element & \textit{Block}\texttt{: a group of }\textit{n}\texttt{ blocks (…) – double-click opens it} \\
\gui{Trace rays} & verfolgt Strahlen entlang der gewählten Strahlengänge, ein Linien-Satz je Strahlengang; mit \gui{Quick irradiance} zusätzlich eine grobe Bestrahlungsstärke als Feld & \texttt{rays reaching the end of the path – }\textit{Pfad}\texttt{: }\textit{ok}\texttt{/}\textit{gesamt}\texttt{, …} \\
\gui{Update drawing} & zeichnet das System des Node-Workflows sofort & \texttt{drawn from the workflow: …} \\
\gui{Workflow from system} & erzeugt den Node-Workflow aus dem aktuellen System (je Strahlengang ein Zweig); übernimmt Änderungen eines früheren Workflows & \texttt{workflow: }\textit{n}\texttt{ blocks, }\textit{m}\texttt{ connections, }\textit{k}\texttt{ branches; kept …} \\
\bottomrule
\end{referenz}

\section{Parameter}\label{sec:ref-rt-parameter}

Die Parameter stehen im Panel \gui{Plugin} in vier Gruppen. Sie werden pro Dokument gespeichert;
\gui{Reset parameters} stellt die Standardwerte her.

\begin{referenz}{@{}>{\raggedright\arraybackslash}p{3.3cm} >{\raggedright\arraybackslash}p{1.5cm} >{\raggedright\arraybackslash}p{1.8cm} >{\raggedright\arraybackslash}p{2.1cm} >{\raggedright\arraybackslash}X@{}}
\toprule
Schlüssel / Beschriftung & Typ & Standard & Bereich, Einheit & Bedeutung \\
\midrule\endhead
\multicolumn{5}{@{}l}{\textbf{Gruppe \gui{System}}} \\
\param{system}\newline\gui{System file} & Datei & leer & Filter \enquote{Raytatouille system (*.rtt.json *.json)} & die Systemdatei für \gui{Load system} \\
\param{catalogs}\newline\gui{Glass catalogues} & Dateiliste & leer & Filter \enquote{Glass catalogues (*.agf *.AGF);;All files (*)} & AGF-Kataloge (\gui{Add…}, \gui{Remove}). Immer zusätzlich verwendet: \datei{*.agf} neben der Systemdatei und in ihrem Unterordner \datei{catalogs}. Dokumente mit der alten Textform (durch \code{;} getrennte Pfade) werden weiterhin gelesen \\
\param{tolerance}\newline\gui{Drawing tolerance} & Zahl & 0{,}0001 & $10^{-7}$ \dots{} 1~mm; 7 Nachkommastellen & größte Abweichung der gezeichneten NURBS von der optischen Fläche (nur die Zeichnung; Strahlen laufen auf den exakten Flächen) \\
\midrule
\multicolumn{5}{@{}l}{\textbf{Gruppe \gui{Trace}}} \\
\param{path}\newline\gui{Paths} & Text & leer & -- & Namen der Strahlengänge, durch Kommas getrennt; leer: alle. Gilt auch für \gui{Camera image} \\
\param{fields}\newline\gui{Field} & Ganzzahl & 0 & $\geq 0$ & 0: alle Feldpunkte; $k$: nur der $k$-te. Gilt auch für \gui{Camera image} \\
\param{wavelengths}\newline\gui{Wavelengths} & Auswahl & \code{all} & \code{all}, \code{reference} & alle Wellenlängen oder nur die Referenzwellenlänge. Gilt auch für \gui{Camera image} \\
\param{rays}\newline\gui{Rays (drawn)} & Ganzzahl & 127 & 1 \dots{} 20\,000 & gezeichnete Strahlen je Feld und Wellenlänge (hexapolare Pupille; die nächste mögliche Zahl 1, 7, 19, 37, \dots{} 127, \dots{} wird verwendet) \\
\param{colour\_by}\newline\gui{Colour by} & Auswahl & \code{wavelength} & \code{wavelength}, \code{field}, \code{path} & Einfärbung der Strahlen \\
\param{irradiance}\newline\gui{Quick irradiance} & Ja/Nein & aus & -- & grobe Bestrahlungsstärke mit \gui{Trace rays}; für ein echtes Kamerabild \gui{Camera image} \\
\param{irradiance\_rays}\newline\gui{Rays (quick irradiance)} & Ganzzahl & 5\,000 & 7 \dots{} 500\,000 & Strahlen je Feld und Wellenlänge der schnellen Bestrahlungsstärke \\
\param{grid}\newline\gui{Quick irradiance grid} & Ganzzahl & 128 & 8 \dots{} 1\,024 & Zellen an der langen Seite der schnellen Bestrahlungsstärke \\
\midrule
\multicolumn{5}{@{}l}{\textbf{Gruppe \gui{Camera image}}} \\
\param{camera\_rays}\newline\gui{Rays per field and wavelength} & Ganzzahl & 200\,000 & 1\,000 \dots{} 50\,000\,000 & viele Strahlen für ein glattes Bild; die Zeit wächst linear \\
\param{camera\_pixels}\newline\gui{Pixels (long side)} & Ganzzahl & 512 & 16 \dots{} 4\,096 & Pixel an der langen Seite; Pixel quadratisch \\
\param{camera\_on\_surface}\newline\gui{Also on the surface in 3D} & Ja/Nein & an & -- & das Bild zusätzlich als Feld auf der Fläche \\
\midrule
\multicolumn{5}{@{}l}{\textbf{Gruppe \gui{Optimise}}} \\
\param{save\_as}\newline\gui{Save as} & Datei zum Spei\-chern & leer & Filter \enquote{Raytatouille system (*.rtt.json)} & Ziel von \gui{Save system}; leer: neben der Systemdatei, Endung \datei{.optimised.rtt.json} statt \datei{.rtt.json} \\
\bottomrule
\end{referenz}

\paragraph{Mögliche Strahlenzahlen.} Die Pupille wird hexapolar mit $r$ Ringen abgetastet:
$N = 1 + 3r(r+1)$ Strahlen, also 7, 19, 37, 61, 91, 127, 169, 217, 271, 331, \dots{} Das Plugin
wählt das $r \geq 1$, dessen $N$ der Eingabe am nächsten liegt. Beispiele: 127 $\to$ 6 Ringe
(127 Strahlen), 5\,000 $\to$ 40 Ringe (4\,921 Strahlen), 200\,000 $\to$ 258 Ringe (200\,467
Strahlen).

\section{Ergebnisse}\label{sec:ref-rt-ergebnisse}

Jede rechnende Aktion legt im Dock \gui{Results} einen neuen Lauf an. Linien und Felder lassen
sich dort einzeln ein- und ausblenden, Bilder erscheinen im Fenster \gui{Image}, Kurven im Dock
\gui{Plot}.

\begin{referenz}{@{}>{\raggedright\arraybackslash}p{2.4cm} >{\raggedright\arraybackslash}p{3.3cm} >{\raggedright\arraybackslash}p{2.1cm} >{\raggedright\arraybackslash}X@{}}
\toprule
Aktion & Lauf (Dock \gui{Results}) & Art & Inhalt \\
\midrule\endhead
\gui{Trace rays} & \texttt{Trace }\textit{Datei} & Linien & \texttt{Rays }\textit{Pfad} je Strahlengang; Werte nach \gui{Colour by}: Wellenlänge in nm (viridis), Feldnummer (plasma) oder Pfadnummer (coolwarm, Bereich 1 bis Zahl der Pfade) \\
 & & Feld auf Fläche & \texttt{Irradiance }\textit{Fläche} (nur mit \gui{Quick irradiance}), Einheit \code{W/mm\textsuperscript{2} (rel.)} \\
\gui{Camera image} & \texttt{Camera image }\textit{Datei} & Bild & \texttt{Camera image }\textit{Fläche} im Fenster \gui{Image}, Achsen in mm (lokale x, y der Fläche), Einheit \code{W/mm\textsuperscript{2} (rel.)}; Beschreibung mit Strahlen, Pfaden, Feldern, Wellenlängen, Pixelzahl und Pixelgröße \\
 & & Feld auf Fläche & \texttt{Irradiance }\textit{Fläche} (mit \gui{Also on the surface in 3D}) \\
\gui{Optimise} & \texttt{Optimisation} & Kurve & \texttt{Merit function} im Dock \gui{Plot}: x \enquote{iteration}, y \enquote{log10 $\varphi$}, wächst live \\
\bottomrule
\end{referenz}

\textit{Datei} ist der Name der Systemdatei bis zum ersten Punkt. Diese Ergebnisse werden nicht im
Dokument gespeichert.

\section{Modellobjekte und Attribute}\label{sec:ref-rt-attribute}

\gui{Load system} (und das Zeichnen aus dem Workflow) legt an:

\begin{itemize}
  \item eine Gruppe \texttt{Raytatouille: }\textit{Name} und darunter Gruppen für Baugruppen und
    Elemente des Systembaums;
  \item je optische Fläche eine NURBS-Fläche, benannt nach der Flächen-\param{id};
  \item je Linsenrand eine Regelfläche \textit{Element}\texttt{ edge} (weitere: \texttt{edge 2}, \dots).
\end{itemize}

\begin{referenz}{@{}>{\raggedright\arraybackslash}p{3.6cm} >{\raggedright\arraybackslash}p{2.6cm} >{\raggedright\arraybackslash}X@{}}
\toprule
Attribut & an & Wert \\
\midrule\endhead
\param{rtt.surface} & jeder Fläche & \param{id} der Fläche im System \\
\param{rtt.location} & jeder Fläche & JSON-Pfad der Fläche im Systembaum \\
\param{rtt.kind} & jeder Fläche & Art der Fläche laut Raytatouille (z.\,B. \code{detector}) \\
\param{rtt.shape} & jeder Fläche & Form der Fläche laut Raytatouille \\
\param{rtt.nurbs\_error\_mm} & jeder Fläche & größte Abweichung der Zeichnung in mm (3 gültige Stellen) \\
\param{rtt.anchor} & erster Fläche & \code{true}: Anker des Systems \\
\param{rtt.revision} & Anker & Nummer des Systemstands im Plugin-Prozess (folgt Rückgängig/Wiederholen) \\
\param{rtt.file} & Anker & Pfad der geladenen Systemdatei \\
\param{rtt.workflow} & Anker & Kennung des Workflow-Stands, wenn das System aus dem Workflow gezeichnet wurde \\
\param{rtt.rim} & Linsenrändern & Name des Elements \\
\param{sim.inactive} & beliebig (von Hand) & wird noch nicht ausgewertet; das Log nennt solche Objekte \\
\bottomrule
\end{referenz}

Objekte mit \param{rtt.surface} oder \param{rtt.rim} und Gruppen, deren Name mit
\texttt{Raytatouille: } beginnt, werden beim nächsten Zeichnen entfernt.

\section{Workflow-Blöcke}\label{sec:ref-rt-bloecke}
\index{Node-Workflow!Raytatouille-Blöcke}

Der Raytracer stellt dem Node-Workflow die Raytatouille-Blocktypen zur Verfügung; berechnet werden
sie im Prozess des Raytracers. Die Typen und ihre Anschlüsse kommen aus dem Paket \code{rtt\_fusion}
(RTT-Fusion); die folgende Liste gibt den Stand der mit diesem Handbuch ausgelieferten Version
wieder und kann sich mit neueren Versionen ändern. Die Buchstaben bei den Anschlüssen nennen die
Darstellungen des Lichts, die ein Anschluss annimmt: \textbf{R} Strahlen, \textbf{H} Felder; die
weiteren (G, S, T, P) beschreibt \autoref{kap:workflow}. Allgemeines zu Blöcken, Anschlüssen und
Zweigen steht in \autoref{kap:workflow} und \autoref{kap:workflow-bloecke}.

\subsection{Blöcke aus einem System}

Diese Blöcke entstehen mit \gui{Workflow from system}.

\begin{referenz}{@{}>{\raggedright\arraybackslash}p{2.4cm} >{\raggedright\arraybackslash}p{4.6cm} >{\raggedright\arraybackslash}X@{}}
\toprule
Typ, Kategorie & Anschlüsse & Parameter \\
\midrule\endhead
\code{rtt.source}\newline Kategorie \code{source}\newline\small Feldpunkte, Wellenlängen, Polarisation eines Systems & \code{out} (aus, Licht \enquote{light}; G, H, P, R, S) & \param{system} (Datei, \enquote{Raytatouille system}); \param{fields} (Text: Feldindizes, leer: alle); \param{wavelengths} (Text: Wellenlängenindizes, leer: alle); Daten: \param{rtt\_system} (Einstellungen des Systems), \param{rtt\_tree} (Baumknoten, die keinem Block gehören), \param{rtt\_paths} (Reihenfolge und Schlüssel der Strahlengänge), \param{rtt\_catalogs} (Glaskatalog-Dateien), \param{rtt\_catalog\_data} (gespeicherte Glaskataloge) \\
\code{rtt.section}\newline Kategorie \code{domain}\newline\small sequentielle Flächen, z.\,B. ein Objektiv & \code{front\_in}, \code{back\_out}, \code{back\_in}, \code{front\_out} (G, H, P, R, S, T) & \param{location} (Text: Knoten im Systembaum); \param{offset} (Punkt, mm: Verschiebung der Elemente des Blocks); \param{drive} (Zahl, mm: Verschiebung entlang der Flächennormalen, gesetzt vom Antrieb eines Phasenschiebers). Kontextmenü: \gui{Section $\to$ group} \\
\code{rtt.mirror}\newline Kategorie \code{domain}\newline\small Spiegel (eine reflektierende Fläche) & \code{in}, \code{out} (G, H, R, S) & \param{surface} (Text); \param{offset} (Punkt, mm); \param{drive} (Zahl, mm) \\
\code{rtt.splitter}\newline Kategorie \code{splitter}\newline\small teilt in Kanäle, vereint kohärent je Gruppe & ein: \code{in}, \code{r\_in}, \code{t\_in}; aus: \code{r} (reflected), \code{t} (transmitted), \code{merged} (combined); jeweils G, H, R, S & \param{surface} (Text); \param{offset} (Punkt, mm); \param{drive} (Zahl, mm) \\
\code{rtt.detector}\newline Kategorie \code{detector}\newline\small Detektor bzw. Kamera: Amplituden innerhalb, Intensitäten zwischen Kohärenzgruppen & \code{in} (G, H, P, R, S); Daten aus: \code{image} & \param{surface} (Text); \param{pixels} (Ganzzahl, \enquote{pixels (long side)}, Standard 512, 16 \dots{} 8\,192); \param{half\_size} (Zahlenliste, mm: Bildfenster $[h_u, h_v]$; leer: die Apertur des Detektors) \\
\bottomrule
\end{referenz}

Die Blöcke eines importierten Systems enthalten außerdem die Daten ihrer Elemente (Parameter
\param{rtt}, \param{runs}); sie erscheinen im Dialog \gui{Parameters…} als Reiter je Element
(Längen in mm, Winkel in Grad bei Schlüsseln mit \code{\_deg}, Wellenlängen in
\si{\micro\metre}).

\subsection{Neue Teile (\code{rtt.part.*})}

Diese Teile fügen Sie mit \gui{Add block} ein. Jedes ist ein echtes Element des Systems, das der
Workflow beschreibt; Lage und Daten geben Sie in \gui{Parameters…} ein. Alle haben nur den
Daten-Parameter \param{rtt\_new} (das Raytatouille-Element: Lage, Flächen, Glas), dessen
Vorgabe hier steht. Im Editor zeigt ein Block je Lichtanschluss einen Punkt; intern hat jeder
Anschluss eine Eingangs- und eine Ausgangsseite (\code{\_in}, \code{\_out}). Alle Lichtanschlüsse
nehmen H und R an.

\begin{referenz}{@{}>{\raggedright\arraybackslash}p{3.3cm} >{\raggedright\arraybackslash}p{3.2cm} >{\raggedright\arraybackslash}X@{}}
\toprule
Typ & Anschlüsse & Vorgabe von \param{rtt\_new} \\
\midrule\endhead
\code{rtt.part.mirror}\newline Raytatouille mirror & \code{front} & \code{mirror} \enquote{Mirror}, Fläche \code{M}, kreisförmige Apertur Radius 12{,}5~mm, Wechselwirkung \code{ideal\_mirror} \\
\code{rtt.part.splitter}\newline Raytatouille beam splitter & \code{front\_a}, \code{front\_b}, \code{back\_a}, \code{back\_b} & \code{thin\_element} \enquote{Beam splitter}, um \SI{45}{\degree} gedreht (\code{rotation\_deg} [45, 0, 0]), Fläche \code{BS}, Radius 12{,}5~mm, \code{ideal\_beam\_splitter} \\
\code{rtt.part.lens}\newline Raytatouille lens & \code{front}, \code{back} & \code{lens} \enquote{Lens}, Material \code{CONST:1.5168}, Flächen \code{S1} (Radius 51{,}68~mm) und \code{S2} (Radius $-51{,}68$~mm, 4~mm dahinter), Aperturradius 12{,}5~mm \\
\code{rtt.part.stop}\newline Raytatouille stop & \code{front}, \code{back} & \code{stop} \enquote{Stop}, Fläche \code{STO}, kreisförmige Apertur Radius 5~mm \\
\code{rtt.part.detector}\newline Raytatouille detector & \code{in}; Daten aus: \code{image} & \code{detector} \enquote{Detector}, Fläche \code{IMG}, rechteckige Apertur $\pm 5$~mm \\
\bottomrule
\end{referenz}

Beim Strahlteiler gilt: Wer das Teil auf der Seite verlässt, auf der das Licht hereinkam, wurde
reflektiert. Ein Spiegel wird über denselben Anschluss verlassen, über den das Licht kam.

\subsection{Übersichtsbild}\label{sec:ref-rt-uebersicht}
\index{Uebersichtsbild@Übersichtsbild!Raytracer}

Für alle Blöcke \code{rtt.*} zeichnet der Raytracer den Reiter \gui{Overview} des Dialogs
\gui{Parameters…} (Modul \datei{rtt\_preview.py}; Abschnitt~\ref{sec:workflow-uebersicht}):

\begin{referenz}{@{}>{\raggedright\arraybackslash}p{3.0cm}X@{}}
\toprule
Block & Inhalt \\
\midrule\endhead
\code{rtt.section}, \code{rtt.part.lens} & Meridionalschnitt (y–z, mm) der Elemente: Linsen als
  geschlossene Kontur je Glas (\code{raytatouille.layout.outlines}, exakte Pfeilhöhen über den
  Aperturen), Legende \enquote{Glas · n = …} bei der Referenzwellenlänge der Quelle, Kittglieder
  je Glas; Blenden als Marken bis zum 1,5-fachen Öffnungsradius mit \enquote{stop $\varnothing$ …}; ideale
  Linsen mit Pfeilspitzen und $f$; optische Achse; Elementnamen (bis zwölf Elemente); sieben
  parallele Strahlen entlang des ersten Laufs (\param{runs}) des Blocks, auf 85\,\% der kleinsten
  Öffnung; Bemaßung Mittendicke bzw.\ erster bis letzter Scheitel. Bei einer einzelnen Linse
  nennen die Hinweise Radien, Konik, Dicken und Durchmesser. \\
\code{rtt.mirror}, \code{rtt.splitter}, \code{rtt.detector}, übrige \code{rtt.part.*} &
  Profil der Fläche in der y–z-Ebene: Spiegel kräftig, Strahlteiler gestrichelt, Detektor blau,
  Blende als Marken; ohne Strahlen. \\
\code{rtt.source} & das ganze System (Baumgerüst \param{rtt\_tree}, die Elemente aller Blöcke, neue
  Teile) ohne Strahlen; Hinweis mit Wellenlängen und Zahl der Feldpunkte. \\
\bottomrule
\end{referenz}

Die Koordinaten sind die der Baugruppe des Blocks (Abschnitte, Spiegel) bzw.\ des Systems (Quelle,
neue Teile). Kompiliert wird mit den Glaskatalogen der Quelle (\param{rtt\_catalog\_data},
\param{rtt\_catalogs}); fehlt ein Glas, wird die Geometrie mit einer festen Brechzahl gezeichnet,
und der Hinweis lautet \enquote{glasses unknown (…) -- geometry only, no rays; add the glass
catalogue in the Raytracer panel}. Flächen, die die y–z-Ebene nicht schneidet, nennt der Hinweis
\enquote{… not in the y–z section}. Ohne das Paket \code{raytatouille} meldet der Reiter
\enquote{the Raytracer's overview needs the Python package raytatouille}.

\section{Glaskataloge}\label{sec:ref-rt-kataloge}

Verwendet werden, ohne Doppelte, in dieser Reihenfolge:
\begin{enumerate}
  \item die Dateien von \gui{Glass catalogues},
  \item \datei{*.agf} im Ordner der Systemdatei,
  \item \datei{*.agf} im Unterordner \datei{catalogs},
  \item im Workflow gespeicherte Kataloge (\param{rtt\_catalog\_data} der Quelle), soweit keine
    Datei sie liefert.
\end{enumerate}
Der Workflow übernimmt die Dateien aus 1.--3. in \param{rtt\_catalogs} und speichert ihren Inhalt
bei \gui{Check workflow} oder \gui{Run workflow} in \param{rtt\_catalog\_data}.

\section{Meldungen im Plugin-Log}\label{sec:ref-rt-log}

\begin{referenz}{@{}>{\raggedright\arraybackslash}p{0.45\linewidth} >{\raggedright\arraybackslash}X@{}}
\toprule
Meldung & Bedeutung \\
\midrule\endhead
\texttt{glass catalogues: …} & verwendete Kataloge; gespeicherte mit \texttt{(stored)} \\
\texttt{catalogue …: …} & ein Katalog ließ sich nicht lesen und wurde übersprungen \\
\texttt{drawn with more than the tolerance …} & Flächen über der Zeichentoleranz, mit Abweichung \\
\texttt{irradiance on … (…): rays not counted – …} & nicht gezählte Strahlen: \texttt{outside the grid}, \texttt{vignetted}, \texttt{not reached} \\
\texttt{irradiance on …: curved surface of revolution – not mapped yet} & keine Bestrahlungsstärke auf gekrümmten Rotationsflächen \\
\texttt{irradiance: … is not drawn} & die letzte Fläche ist nicht gezeichnet \\
\texttt{… the last event changes the weight – computed with recorded paths and fewer rays (… pupil rings)} & reflektierende/absorbierende letzte Fläche: höchstens 80 Ringe \\
\texttt{path …: …} & ein Strahlengang konnte nicht verfolgt werden (die übrigen schon) \\
\texttt{marked inactive in CAD but still traced …} & Objekte mit \param{sim.inactive}, die trotzdem gerechnet werden \\
\texttt{drawing not updated: …}, \texttt{fields not updated: …} & \gui{Follow the workflow} konnte nach einer Workflow-Änderung nicht zeichnen bzw. die Felder der Blöcke nicht beschreiben \\
\bottomrule
\end{referenz}
